\chapter{Technical details for mean-preserving and mean-tilted intervals}
\label{app:mti-technical}

This appendix preserves the detailed derivations, algorithms, and
secondary diagnostics supporting the corresponding research chapter.
The chapter retains the inferential target, central model, principal
evidence, interpretation, and limitations needed for a continuous read.

% migration-id: MIG-CH05-01
\section{Score derivation and consequences}
\label{rqr:sec:supp-population-target}

The derivation starts from the endpoint scores because those scores determine
both content and placement. Fix a covariate value \(x\), a desired content
\(c\in(0,1)\), and endpoints
\(a=L(x)<U(x)=b\). The pointwise loss is
\begin{equation}
\ell_c(a,b;y)
=
\begin{cases}
(1-c)(y-a)(b-y), & a<y<b,\\
c(y-a)(y-b), & y\le a\ \text{or}\ y\ge b.
\end{cases}
\label{rqr:eq:supp-mpi-loss-piecewise}
\end{equation}
Equivalently, with
\(e(a,b;y)=(y-a)(y-b)\) and
\(\rho_c(r)=r\{c-\ind(r<0)\}\),
\begin{equation}
\ell_c(a,b;y)=\rho_c\{e(a,b;y)\}.
\label{rqr:eq:supp-mpi-product-check}
\end{equation}
The check function is applied to the product residual. Its index \(c\) denotes
the desired interval content rather than a prescribed response-quantile level
for either endpoint or a generalized-Bayes learning rate. The loss is symmetric
in its two roots; ordering is used only for the endpoint interpretation.

Define \(I_{a,b}(Y)=\ind(a<Y<b)\). Away from the two roots, the derivative of
the check loss with respect to its scalar argument is
\(c-I_{a,b}(Y)\). Assume
\(\E(Y^2\mid X=x)<\infty\), the conditional risk is differentiable at an
interior stationary point, differentiation and conditional expectation may be
interchanged, and the conditional distribution places no mass at \(a\) or
\(b\). Since
\[
\frac{\partial}{\partial a}(Y-a)(Y-b)=b-Y,
\qquad
\frac{\partial}{\partial b}(Y-a)(Y-b)=a-Y,
\]
the pointwise first-order conditions are
\begin{align}
\E\left[\{c-I_{a,b}(Y)\}(b-Y)\mid X=x\right]&=0,
\label{rqr:eq:pop-foc-a}\\
\E\left[\{c-I_{a,b}(Y)\}(a-Y)\mid X=x\right]&=0.
\label{rqr:eq:pop-foc-b}
\end{align}
Subtracting \eqref{rqr:eq:pop-foc-b} from \eqref{rqr:eq:pop-foc-a} and using
\(b-a>0\) gives
\[
\Pr\{a<Y<b\mid X=x\}=c.
\]
Substitution into either first-order condition then gives
\[
\E\{Y\ind(a<Y<b)\mid X=x\}=c\E(Y\mid X=x).
\]

The second equation is equivalently
\[
\E(Y\mid a<Y<b,X=x)=\E(Y\mid X=x)
\]
and
\(\operatorname{Cov}\{Y,I_{a,b}(Y)\mid X=x\}=0\). This covariance statement is
compatible with dependence, asymmetry, and unequal-tailed endpoints. The
coverage implication is established by \citet{rqr-PouplinEtAl2024RQR}; the
conditional first-moment implication is a further consequence of their
first-order system.

Let \(\mu_x=\E(Y\mid X=x)\). Since
\(\E(Y-\mu_x\mid X=x)=0\), the first-moment equation is also equivalent to
\begin{equation}
\E[\{\mu_x-Y\}\ind(Y<a)\mid X=x]
=
\E[\{Y-\mu_x\}\ind(Y>b)\mid X=x].
\label{rqr:eq:rqr-tail-moment-balance}
\end{equation}
Thus MPI balances omitted tail probability multiplied by average
tail displacement from the conditional mean. The two omitted tail probabilities
may differ.

The proof of Theorem~\ref{rqr:thm:quantile-window} follows directly from the
probability-window representation and is included here to keep the assumptions,
statement, and argument together.
\begin{proof}
Modulo conditional probability-zero portions outside the support, every
contiguous interval having content \(c\) has the displayed
quantile-window form. If \(v>u\), then
\[
M_{c,x}(v)-M_{c,x}(u)
=\frac{1}{c}\int_0^c
\{Q_x(v+s)-Q_x(u+s)\}\,ds>0.
\]
The average of the lowest conditional \(c\)-fraction is below \(\mu_x\), and
the average of the highest conditional \(c\)-fraction is above \(\mu_x\), for
a nondegenerate distribution with finite mean. Continuity and strict
monotonicity therefore give existence and uniqueness. The first-moment
equation identifies the stated endpoints.
\end{proof}

If \(F_x\) is symmetric about \(\mu_x\), then
\(Q_x(1-v)=2\mu_x-Q_x(v)\). The window beginning at
\((1-c)/2\) has retained mean \(\mu_x\), so uniqueness gives the equal-tailed
interval. Conversely, coincidence at one coverage level does not establish
global symmetry.

The distribution conditional on interval membership and its complementary
tail distribution both have mean \(\mu_x\). This is a first-moment statement,
not independence between \(Y\) and interval membership. It also clarifies why
the endpoint midpoint need not equal \(\mu_x\) under asymmetry.

\begin{proposition}[Mean-preserving contraction]
\label{rqr:prop:supp-convex-order}
Let \(Y_{\mathrm{in}}\) and \(Y_{\mathrm{out}}\) have the conditional
distributions of \(Y\) given \(a<Y<b\) and \(Y\notin(a,b)\), respectively,
at a regular MPI target. For every convex function \(\varphi\) for
which the required expectations exist,
\[
\E\{\varphi(Y_{\mathrm{in}})\}
\le
\E\{\varphi(Y)\}
\le
\E\{\varphi(Y_{\mathrm{out}})\}.
\]
Consequently,
\[
Y_{\mathrm{in}}\preceq_{\mathrm{cx}}Y
\preceq_{\mathrm{cx}}Y_{\mathrm{out}}.
\]
If second moments exist, the same ordering holds for the corresponding
variances.
\end{proposition}

\begin{proof}
Let \(s_\varphi\) be the affine secant through
\(\{a,\varphi(a)\}\) and \(\{b,\varphi(b)\}\). Convexity gives
\(\varphi(y)\le s_\varphi(y)\) on \([a,b]\) and
\(\varphi(y)\ge s_\varphi(y)\) outside that interval. Both conditional
distributions have mean \(\mu_x\), so
\[
\E\{\varphi(Y_{\mathrm{in}})\}
\le s_\varphi(\mu_x)
\le \E\{\varphi(Y_{\mathrm{out}})\}.
\]
The full conditional distribution is their mixture with weights \(c\) and
\(1-c\), which places its convex expectation between the two. Taking
\(\varphi(y)=(y-\mu_x)^2\) gives the variance statement.
\end{proof}

\section{Interval Paths and Tail Sensitivity}

Retain the quantile-window characterization and assume, for this derivative
calculation, connected support and positive continuous density at the relevant
endpoint quantiles. Write \(u_c\) for the MPI index, and set
\(L_c=Q(u_c)\) and \(U_c=Q(u_c+c)\). Differentiating
\[
\int_{u_c}^{u_c+c}\{Q(v)-\mu\}\,dv=0
\]
with respect to \(c\), at regular interior values, gives
\[
\frac{du_c}{dc}
=-\frac{U_c-\mu}{U_c-L_c}<0,
\qquad
\frac{d(u_c+c)}{dc}
=\frac{\mu-L_c}{U_c-L_c}>0.
\]
The MPI population intervals are therefore nested as content
increases, with analogous one-sided statements at isolated nonsmooth support
transitions. If \(p_\mu=F(\mu)\) satisfies \(Q(p_\mu)=\mu\) and \(Q\) is
continuous at \(p_\mu\), then \(Q(u_c)\le\mu\le Q(u_c+c)\), so
\(p_\mu\in[u_c,u_c+c]\). Since this probability window has length \(c\),
\[
L_c\longrightarrow\mu,\qquad U_c\longrightarrow\mu
\quad\text{as }c\downarrow0.
\]
This zero-content limit needs the local support condition above. It is false
under the quantile-window assumptions alone: for the symmetric distribution
with density \(1/2\) on \([-2,-1]\cup[1,2]\), \(\mu=0\) and the MPI
window has \(u_c=(1-c)/2\), hence \(L_c=-1-c\) and \(U_c=1+c\), so the
zero-content limit straddles the support gap rather than contracting to
\(\mu\). The corrected zero-content limit is a property of the population
interval path; the
midpoint at a fixed positive \(c\) is not generally the mean, and independent
fits at several content levels do not by themselves define joint
generalized-posterior draws of that path.

For \(Y^\star=r+sY\), \(s>0\), the MPI endpoints transform as
\[
L_c^\star=r+sL_c,\qquad U_c^\star=r+sU_c.
\]
For the tilted family, \(\delta^\star=s\delta\) and
\[
\ell^\star_{c,\delta^\star}
(r+sa,r+sb;r+sy)
=s^2\ell_{c,\delta}(a,b;y).
\]
Thus the population target is positive-affine equivariant. A corresponding
generalized posterior uses the transformed prior and
\(\omega_{\mathrm R}^\star=\omega_{\mathrm R}/s^2\). Conversely, if the
original response is standardized as \(\widetilde Y=(Y-r)/s\), use
\(\widetilde\delta=\delta/s\) and
\(\widetilde\omega_{\mathrm R}=\omega_{\mathrm R}s^2\).
Nonlinear monotone transformations generally change the mean-preserving
functional rather than merely re-expressing it.

Finally, with fixed endpoints,
\(\rho_c\{(y-a)(y-b)\}=O(y^2)\) as \(|y|\to\infty\), and the endpoint scores
grow linearly in \(|y|\). MPI is distribution-free in the sense that it does
not posit a parametric response density. Its retained-mean target still depends
on second moments, so the finite-second-moment assumption and response-scale
choice are scientifically consequential.

Endpoint atoms require directional conditions. Write
\(m=(a+b)/2\), \(h=(b-a)/2>0\), so
\(e=(Y-m)^2-h^2\). The left and right derivatives in \(h\) are
\[
D_-R_x(m,h)=2h\{\Pr(a<Y<b\mid x)-c\},
\]
and
\[
D_+R_x(m,h)=2h\{\Pr(a\le Y\le b\mid x)-c\}.
\]
At a local minimum, \(D_-R_x\le0\le D_+R_x\), and therefore
\[
\Pr(a<Y<b\mid X=x)\le c\le\Pr(a\le Y\le b\mid X=x).
\]
Without a separate joint subgradient and endpoint-mass allocation analysis,
these directional conditions do not imply the exact retained-mean equality
derived above under endpoint continuity.
If \(a=b=m\), then \(e=(Y-m)^2\) and the risk is
\(c\E\{(Y-m)^2\mid X=x\}\); the coincident-root target is the conditional
mean when the second moment exists.

These implications are conditional, pointwise, and assumption dependent. They
do not follow by treating estimated data-dependent endpoints as fixed. For
linear roots \(\eta_k(X)=X^\top\beta_k\), the corresponding population score
conditions are instead
\begin{align}
\E\left[X\{c-I_{\eta_1,\eta_2}(Y)\}(\eta_2-Y)\right]&=0,\\
\E\left[X\{c-I_{\eta_1,\eta_2}(Y)\}(\eta_1-Y)\right]&=0,
\end{align}
where \(I_{\eta_1,\eta_2}(Y)\) indicates that \(Y\) lies between the
unordered roots. These design-weighted projection equations need not imply
conditional coverage for every \(x\).

The finite-sample unbiasedness calculation in
\citet{rqr-PouplinEtAl2024RQR} treats fitted, data-dependent endpoints as fixed,
and its zero-gradient argument does not cover nonsmooth optima when an
observation coincides with a fitted root. We therefore do not use that
calculation as a finite-sample guarantee; such a result requires a separate
empirical-process analysis.

% migration-id: MIG-CH05-02
\section{Identification proofs and restricted-score equations}
\label{rqr:sec:supp-mean-tilt}

The endpoint-score derivation below proves
Proposition~\ref{rqr:prop:mean-tilt} and records the corresponding equations for
restricted root classes. The loss is invariant to exchanging the raw root
labels. Its tilt has response units and changes target placement rather than
the content score.
For restricted linear roots \(\eta_k(X)=X^\top\beta_k\) and an externally
fixed \(\delta(X)\), the population score equations are
\begin{align}
\E\!\left[
X\left\{\{c-I_{\eta_1,\eta_2}(Y)\}(\eta_2-Y)-c\delta(X)\right\}
\right]&=0,\\
\E\!\left[
X\left\{\{c-I_{\eta_1,\eta_2}(Y)\}(\eta_1-Y)-c\delta(X)\right\}
\right]&=0.
\end{align}
They are design-weighted projection equations. Pointwise conditional content
and retained-mean identities require a sufficiently rich predictor class and
the unrestricted conditions below.

\subsection{Fixed-content target and proof}

The following argument establishes the fixed-content and retained-mean
identities in Proposition~\ref{rqr:prop:mean-tilt}; the subsequent profiling
lemma strengthens stationary-point identification to the global result in
Theorem~\ref{rqr:thm:profiled-mean-tilt}.
\begin{proof}
Let \(I_{m,h}(Y)=\ind(|Y-m|<h)\). Away from endpoint atoms,
\begin{align}
\frac{\partial}{\partial h}\E\{\ell_{c,\delta}(m,h;Y)\}
&=2h\{\Pr(I_{m,h}=1)-c\},\\
\frac{\partial}{\partial m}\E\{\ell_{c,\delta}(m,h;Y)\}
&=2\E[(m-Y)\{c-I_{m,h}(Y)\}]-2c\delta.
\end{align}
The first equation gives content \(c\). Substitution into the second gives
\(\E\{YI_{m,h}(Y)\}=c(\mu+\delta)\), which is the retained-mean equation.
Theorem~\ref{rqr:thm:quantile-window} and strict monotonicity of \(M_c\)
give the quantile-window statement.
\end{proof}

\begin{lemma}[Profile radius and envelope derivative]
\label{rqr:lem:supp-profile-radius}
Suppose that the support closure of \(F\) is an interval
\([\ell,r]\), where either endpoint may be infinite, \(F\) is continuous and
strictly increasing and absolutely continuous between its support endpoints,
its density is positive and continuous on the support interior, and
\(\E(Y^2)<\infty\). Extend \(F\) by zero and one beyond finite support
endpoints. For every \(m\in\R\), there is a unique \(h_c(m)>0\) satisfying
\[
\Pr\{|Y-m|<h_c(m)\}=c.
\]
The map \(m\mapsto h_c(m)\) is continuous and is differentiable whenever its
two roots are interior to the support. Define
\[
\overline R_{c,\delta}(m)
=
\E\{\ell_{c,\delta}(m,h_c(m);Y)\}.
\]
Also let \(u_c(m)=F\{m-h_c(m)\}\). Then
\begin{equation}
\overline R_{c,\delta}'(m)
=2c\left[M_c\{u_c(m)\}-\mu-\delta\right],
\label{rqr:eq:supp-profiled-risk-derivative}
\end{equation}
where both roots are interior, with the corresponding one-sided derivatives
when a root crosses a finite support endpoint.
\end{lemma}

\begin{proof}
For fixed \(m\), the map
\[
H_m(h)=F(m+h)-F(m-h),\qquad h\ge0,
\]
is continuous, begins at zero, and converges to one. Interval support and
strict increase imply that \(H_m\) passes through every value in \((0,1)\)
strictly: any plateau can occur only before the expanding interval reaches
the support or after it contains the full support. Thus \(H_m(h)=c\) has one
solution. Continuity of this inverse in \(m\) follows by the same strict
crossing argument. When both roots are interior, the implicit-function
theorem applies because
\[
\frac{\partial H_m(h)}{\partial h}
=f(m+h)+f(m-h)>0.
\]

For fixed \(m\), write \(Z_m=(Y-m)^2\) and \(q=h^2\). The subgradient interval
of \(q\mapsto\E\{\rho_c(Z_m-q)\}\) is
\[
\left[\Pr(Z_m<q)-c,\ \Pr(Z_m\le q)-c\right].
\]
Continuity makes \(q=h_c(m)^2\) its unique minimizer. On a compact
neighborhood of \(m\), the partial midpoint score is dominated by a constant
multiple of \(1+|Y|\), which is integrable. Differentiation under the
expectation is therefore valid away from the probability-zero endpoints.
The standard two-sided profiling inequalities, evaluated once at
\(h_c(m)\) and once at \(h_c(m+s)\), together with continuity of \(h_c\),
give the envelope derivative
\begin{align*}
\overline R_{c,\delta}'(m)
&=
2\E\!\left[(m-Y)
\{c-\ind(m-h_c(m)<Y<m+h_c(m))\}\right]-2c\delta\\
&=
2c\left[M_c\{u_c(m)\}-\mu-\delta\right].
\end{align*}
The same inequalities give one-sided derivatives at a finite-support
transition.
\end{proof}

\begin{proof}[Proof of Theorem~\ref{rqr:thm:profiled-mean-tilt}]
Write \(a_c(m)=m-h_c(m)\) and \(b_c(m)=m+h_c(m)\). Continuity gives
\[
F\{b_c(m)\}-F\{a_c(m)\}=c,
\]
so \(F\{b_c(m)\}=u_c(m)+c\). When \(0<u_c(m)<1-c\),
\[
a_c(m)=Q\{u_c(m)\},\qquad
b_c(m)=Q\{u_c(m)+c\},
\]
and therefore
\[
m=m_c\{u_c(m)\}
:=\frac{Q\{u_c(m)\}+Q\{u_c(m)+c\}}{2}.
\]
The map \(m_c(u)\) is strictly increasing. If a root extends beyond a finite
support endpoint, \(u_c(m)\) is clamped at \(0\) or \(1-c\); hence \(u_c\) is
globally nondecreasing.

Lemma~\ref{rqr:lem:supp-profile-radius} gives
\eqref{rqr:eq:supp-profiled-risk-derivative}. Strict monotonicity of \(Q\) also
gives, for \(v>u\),
\[
M_c(v)-M_c(u)
=\frac{1}{c}\int_0^c\{Q(v+s)-Q(u+s)\}\,ds>0.
\]
For an interior admissible tilt, the profiled derivative is therefore negative
before the unique solution of \(M_c(u)=\mu+\delta\) and positive afterward.
The profiled risk has a unique global midpoint minimizer, and the preceding
fixed-\(m\) argument gives its unique positive half-width. This proves joint
global identification of the ordered roots.
\end{proof}

The admissible population tilt range is
\[
M_c(0)-\mu\le\delta\le M_c(1-c)-\mu,
\]
and both endpoint values are finite under the assumed second moment. Interior
tilts correspond one-to-one with finite-root quantile windows. At
\(\delta_- = M_c(0)-\mu\), the lower boundary member is
\([Q(0),Q(c)]\). If \(Q(0)=-\infty\), it is a semi-infinite limiting interval
with a finite population-risk infimum. If \(Q(0)\) is finite, every pair
\((a,Q(c))\) with
\(a\le Q(0)\) has the same population loss, so a support constraint is needed
to choose a unique inactive lower root. The upper boundary has the symmetric
qualification. For
\(\delta<\delta_-\) or \(\delta>\delta_+\), the unrestricted population risk
is unbounded below along the corresponding escaping-root direction and has no
finite global target.

MPI and equal-tailed recovery use
\[
\delta_{\mathrm{MPI}}=0,\qquad
\delta_{\mathrm{ET}}=M_c\{(1-c)/2\}-\mu,
\]
whereas shortest-contiguous recovery is generally set-valued. Define
\[
\Delta_{\mathrm{SH}}
=
\{M_c(u)-\mu:u\in S_{\mathrm{SH}}\}.
\]
If \(S_{\mathrm{SH}}\) is a singleton, its element and tilt are denoted
by \(u_{\mathrm{SH}}\) and \(\delta_{\mathrm{SH}}\). These identities establish
population representation. Finite-sample recovery-tilt estimation requires an
external selection protocol.

Implicit differentiation yields
\[
\frac{du_\delta}{d\delta}
=\frac{c}{W_c(u_\delta)}
\]
and
\[
\frac{dW_c(u_\delta)}{d\delta}
=\frac{c\,W_c'(u_\delta)}{W_c(u_\delta)}.
\]
Hence profiling population width along the interior admissible tilt path
identifies any interior shortest-contiguous member. Boundary members use the
one-sided closure described above, and the statement does not extend
the target to disconnected highest-density regions.

\section{Cornish--Fisher Tilt Approximations}
\label{rqr:sec:supp-cf-tilt}

The recovery identities also give a simple initialization approximation. Let
\(Z=(Y-\mu)/\sigma\) have small standardized skewness \(\gamma_1\), and let
\(z_p=\Phi^{-1}(p)\). The formal first-order Cornish--Fisher expansion
\citep{rqr-CornishFisher1937}
\[
Q_Z(p)
 \approx
z_p+\frac{\gamma_1}{6}(z_p^2-1)
\]
approximates the quantile function of \(Z\) around the standard Normal law.
For content \(c\), define
\[
q_c=\Phi^{-1}\{(1+c)/2\},\qquad
K_c=q_c\phi(q_c)/c.
\]
Integrating the expansion over the equal-tailed probability window gives
\[
d_{\mathrm{ET}}^{\mathrm{CF}}
\approx
-\frac{\gamma_1}{3}K_c,
\qquad d=\delta/\sigma.
\]
Applying the same expansion to the shortest-window first-order condition
\(W_c'(u)=0\) and then evaluating the retained-window mean gives
\[
d_{\mathrm{SH}}^{\mathrm{CF}}
\approx
-\gamma_1K_c.
\]
Thus the shortest-oriented first-order tilt is three times the equal-tailed
one on the standardized scale. The sign convention is useful for checking
numerical calculations: right skewness gives negative recovery tilts, while left
skewness gives positive recovery tilts.

These displays are used as first-order approximations, not as a universal remainder
theorem. An \(O(\epsilon^2)\) statement is defensible only for an explicitly
indexed near-Normal family \(F_\epsilon\) whose quantile function and
derivative admit a uniform expansion
\[
Q_\epsilon(p)
=z_p+\epsilon\kappa_3(z_p^2-1)/6+r_\epsilon(p),
\qquad
\sup_p\{|r_\epsilon(p)|+|\partial_pr_\epsilon(p)|\}=O(\epsilon^2),
\]
on a neighborhood of the central window, with
\(\gamma_1(F_\epsilon)=\epsilon\kappa_3+O(\epsilon^2)\). Small skewness by
itself does not control higher cumulants or the quantile remainder.

For a known population law these formulas use the true \(\gamma_1\) and
\(\sigma\). A finite-sample plug-in version may replace them by the adjusted
Fisher--Pearson skewness \(\widehat\gamma_1\) and training standard deviation
\(\widehat\sigma\):
\[
\widehat d_{\mathrm{SH}}^{\mathrm{CF}}
=
-\widehat\gamma_1K_c,\qquad
\widehat d_{\mathrm{ET}}^{\mathrm{CF}}
=
\widehat d_{\mathrm{SH}}^{\mathrm{CF}}/3,\qquad
\widehat\delta=\widehat\sigma\,\widehat d.
\]
The plug-in quantities are fixed tilt choices for initialization and external
validation. A posterior treatment of tilt, or reuse of
the MPI hierarchical-scale update at nonzero tilt, would require a separate
model. Strong skewness, finite-support boundary solutions, and heavy-tailed
cases can move the first-order shortest approximation outside the admissible
finite-root tilt range. For this reason, the numerical summaries include adjusted
skewness, excess kurtosis, standardized sixth moment, probability-window
clipping, boundary proximity, and bootstrap sensitivity summaries. The
deterministic Cornish--Fisher figure and population comparison table
illustrate both near-Normal accuracy and material departures under stronger
skewness or support-boundary geometry. These calculations assess the
approximation, while empirical validation and held-out selection rules require
separate study. Table~\ref{rqr:tab:mean-tilt-cf-mini-study} reports the
population values used in the comparison.

The tilt indexes the estimand and is initially treated like a fixed quantile
level. A scalar tilt generally defines only a restricted global target in a
regression model; pointwise equal-tailed or shortest recovery generally
requires \(\delta(x)\). Sampling \(\delta\) as an ordinary unknown would mix
different interval functionals. Data-driven tilt selection therefore requires
a separate selection rule based on independent test data. Under
positive affine transformation \(Y^\star=r+sY\), the corresponding target uses
\(\delta^\star=s\delta\), as stated in
Section~\ref{rqr:sec:supp-population-target}.

These population equations do not establish finite-sample propriety under
arbitrary root priors. For a fixed learning rate and a finite-dimensional
Gaussian root prior, the MPI loss is nonnegative and the tilt is
linear in the roots. The resulting generalized posterior is proper because a
Gaussian density has finite exponential moments in every linear direction. This covers
the initial ridge fixed-design model and finite-horizon Gaussian state model.
Nonzero tilt under heavier-tailed shrinkage priors requires a separate
coercivity or propriety result.

% migration-id: MIG-CH05-03
\subsection{Protocol, calibration, and feasibility details}
\label{rqr:sec:supp-tolerance-validation-protocol}

The primary tolerance validation is restricted to iid continuous univariate
samples. This is the setting in which the scan-calibrated empirical procedure,
the stability-restricted MTI-ECM profile rule, Young--Mathew interpolation, and
Wilks-style full-range intervals can be compared on the same target. TCSP
is the primary empirically calibrated procedure
because it selects the shortest closed order-statistic window at the calibrated
retained count. The reported TCSP calculations use adaptive Clopper--Pearson scan
calibration directly. The MTI-ECM comparator
profiles fixed-target endpoint fits over fitted content and retained-mean tilt,
then applies a direct Dirichlet-process fixed-interval content probability
threshold and a width-stability restriction. Young--Mathew and Wilks are included as classical order-statistic
comparators. Regression and dynamic tolerance calibration require additional
results because the fitted intervals vary with covariates or time.

The simulation design fixes tolerance confidence \(1-\alpha=0.95\) and uses
explicit feasible design cells:
\[
  \begin{aligned}
  (n,c)\in\{&
  (50,0.90),(100,0.90),(100,0.95),\\
  &(500,0.90),(500,0.95),(500,0.99),\\
  &(1000,0.90),(1000,0.95),(1000,0.99)\}.
  \end{aligned}
\]
The distributional panel contains two symmetric reference distributions and six
asymmetric simulation distributions. The symmetric reference distributions are
Gaussian and Student \(t_3\). The asymmetric simulation distributions are
exponential, asymmetric Laplace with \(\tau=0.10\), two-piece Normal with scale
ratio \(1{:}12\), Beta\((1,8)\), Gamma\((0.5,1)\), and log-normal with
log-scale standard deviation \(1.25\). All laws are centered and scaled before
sampling. Each cell uses paired random samples across methods. The
reproducibility materials include the code version, simulation settings,
package versions, random seeds, calibration files, and numerical summaries.
The comparison combines the completed TCSP, Young--Mathew, and Wilks
confirmatory run with the stability-restricted adaptive MTI-ECM profile rule applied to the
completed independent validation run. The reported rows contribute 288,000
replication-level method evaluations over the 72 distribution-level cells used in the main
comparison. The broader MTI-ECM refinement run is retained as calibration
evidence, with only the selected rule evaluated under the MTI-ECM label in the
independent validation run. All simulations completed, and all four reported methods produced
an interval in every replication. The chapter tables and figures are
computed directly from these replication-level simulation outputs and the
corresponding scan-calibration outputs. The primary table summarizes results
across distributions within each \((n,c)\) cell. Distribution-level
content-attainment percentages, interval-production indicators, infeasibility
indicators, and 2.5\%--97.5\% width ranges remain in the external reproducibility
record rather than being repeated as chapter displays.

The reported MTI-ECM comparator uses a prespecified adaptive
calibration rule indexed by \((n,c,1-\alpha)\). The rule chooses a
direct Dirichlet-process content-probability check from a separate calibration run,
requires width stability relative to TCSP within the same distribution-level
cells, and then applies that choice unchanged to fresh paired replications in
the independent validation
run. Table~\ref{rqr:tab:supp-mti-ecm-policy} records the selected content-probability check, calibration
lower bound, independent validation range, and width span for each cell.

For each simulation law and content level, the shortest population
content-\(c\) interval under the known centered and scaled distribution was
used as a lower width benchmark. It is a population reference, not a
sample-based tolerance procedure.

These results are empirical repeated-sampling evidence for the selected
profile rule. They do not supply the open finite-sample proof for the reported
closed-window TCSP action.

\begingroup
\scriptsize
\setlength{\tabcolsep}{3pt}
\begin{longtable}{@{}>{\raggedright\arraybackslash}p{0.16\textwidth}>{\centering\arraybackslash}p{0.12\textwidth}>{\centering\arraybackslash}p{0.17\textwidth}>{\centering\arraybackslash}p{0.16\textwidth}>{\centering\arraybackslash}p{0.12\textwidth}>{\centering\arraybackslash}p{0.19\textwidth}@{}}
\caption{\textbf{Selected adaptive MTI-ECM calibration rule.} The direct-DP content-probability check is fixed by a separate calibration run within each sample-size/content cell before the independent validation summaries are computed. Validation attainment is reported as the range across the eight simulation distributions.}\label{rqr:tab:supp-mti-ecm-policy}\\
\toprule
Cell & DP content level & Calibration lower probability & Validation range (\%) & Cells below 95\% & Width 95\% span
\\
\midrule
\endfirsthead
\toprule
Cell & DP content level & Calibration lower probability & Validation range (\%) & Cells below 95\% & Width 95\% span
\\
\midrule
\endhead
\(n=50,c=0.90\) & 0.995 & 0.966 & 95.8--97.8 & 0 & 1.53--13.76 \\
\(n=100,c=0.90\) & 0.995 & 0.986 & 98.5--99.7 & 0 & 1.5--5.75 \\
\(n=100,c=0.95\) & 0.995 & 0.962 & 95.1--97.5 & 0 & 2.28--17.68 \\
\(n=500,c=0.90\) & 0.995 & 0.983 & 97.3--99.7 & 0 & 1.25--3.86 \\
\(n=500,c=0.95\) & 0.995 & 0.983 & 97.0--99.2 & 0 & 1.99--5.45 \\
\(n=500,c=0.99\) & 0.995 & 0.954 & 95.2--97.4 & 0 & 4.54--28.5 \\
\(n=1000,c=0.90\) & 0.995 & 0.985 & 97.7--99.5 & 0 & 1.24--3.67 \\
\(n=1000,c=0.95\) & 0.995 & 0.987 & 96.7--99.9 & 0 & 1.96--4.73 \\
\(n=1000,c=0.99\) & 0.996 & 0.984 & 97.9--99.4 & 0 & 4.62--13.31 \\
\bottomrule
\end{longtable}
\endgroup


The TCSP calibration chooses the retained count \(k\) before seeing the
replicated samples. Table~\ref{rqr:tab:supp-tcsp-scan-calibration} records the
retained counts, retained fractions, content buffers, and lower confidence
bounds used in the simulation design. Adaptive Clopper--Pearson calibration
selects the retained count for each feasible sample-size/content cell before
any distribution-specific replications are evaluated. When \(k=n\), TCSP is forced to
use the full sample range and therefore coincides with Wilks as a reported
interval. This occurs for \((n,c)=(50,0.90),(100,0.95)\), and
\((500,0.99)\).

\begin{table}[H]
\centering
\caption{\textbf{TCSP scan calibration for the simulation design.} The retained
count \(k\) is selected by adaptive Clopper--Pearson Monte Carlo scan
calibration. The content buffer is \(k/n-c\).}
\label{rqr:tab:supp-tcsp-scan-calibration}
\TableStyle
\begin{tabular}{@{}rrrrrr@{}}
\toprule
\(n\) & \(c\) & Retained count \(k\) & \(k/n\) & Buffer \(k/n-c\) & Lower confidence bound\\
\midrule
50 & 0.90 & 50 & 1.000 & 0.100 & 0.966 \\
100 & 0.90 & 98 & 0.980 & 0.080 & 0.979 \\
100 & 0.95 & 100 & 1.000 & 0.050 & 0.963 \\
500 & 0.90 & 467 & 0.934 & 0.034 & 0.950 \\
500 & 0.95 & 487 & 0.974 & 0.024 & 0.951 \\
500 & 0.99 & 500 & 1.000 & 0.010 & 0.960 \\
1000 & 0.90 & 925 & 0.925 & 0.025 & 0.960 \\
1000 & 0.95 & 968 & 0.968 & 0.018 & 0.958 \\
1000 & 0.99 & 998 & 0.998 & 0.008 & 0.975 \\
\bottomrule
\end{tabular}

\end{table}

Small-sample feasibility is governed by the full-range two-sided order-statistic
coverage probability
\[
  \Pr\{Y_{(1)}\le Y_{\mathrm{new}}\le Y_{(n)}
       \text{ for at least content } c\}
  =
  1-n(1-c)c^{n-1}-c^n .
\]
The smallest \(n\) satisfying this inequality depends on both content and
tolerance confidence. At confidence \(0.90\), the minimum sample sizes for
contents \(0.90\), \(0.95\), and \(0.99\) are 38, 77, and 388; at confidence
\(0.95\), they are 46, 93, and 473. These thresholds determined the feasible
simulation cells.

This distinction matters for interpreting the small sample sizes in
\citet{rqr-PourmohamadRichardsonSanso2026CalibratedBNP}. The values
\((38,77,388)\) are appropriate at \(90\%\) tolerance confidence. A
\(95\%\)-confidence design needs the larger thresholds \((46,93,473)\). The
simulation design therefore lists the feasible cells explicitly instead of
forming a Cartesian product of sample sizes and contents. This prevents the
comparison from treating structural infeasibility as a method failure.
The Young--Mathew comparator is evaluated through the \texttt{tolerance}
package's \texttt{nptol.int(..., method = "YM")} function. The completed
run used \texttt{tolerance} version 3.0.0, and the package version is included
in the method reproducibility record.

% migration-id: MIG-CH05-04
\section{Data, secondary response, and sensitivity analyses}
\label{rqr:sec:supp-pharma-application}

The pharmaceutical application uses the public manufacturing data described by
\citet{rqr-ZagarMihelic2022PharmaManufacturing}. The analysis file is
\texttt{Laboratory.csv}, version 1 on Figshare
\citep{rqr-ZagarMihelic2021LaboratoryCsv}. The downloaded file contains 1005
batches and 55 recorded fields. The analysis script verifies the byte count and
MD5 digest before restricting the analysis to product code 23. This restriction
leaves 187 batches with recorded strength \(20\)M and size 960000.
The analysis is descriptive and uses procedures fixed before examining these
responses. It illustrates endpoint placement on observed manufacturing batches;
the tolerance-confidence evidence comes from the controlled
simulation study.

The tensile-strength response has 82 distinct values among 187 batches, largest
tie 11, median 1.341, standard deviation 0.124, and skewness 0.44. Weight RSD
has 83 distinct values, largest tie 6, median 1.030, standard deviation 0.467,
and skewness 3.00. Film-coated-tablet hardness was excluded because it has only
32 distinct values and a largest tie of 17 on the recorded scale (median
70.560, standard deviation 6.459, skewness 0.48). Weight RSD is retained as a
strongly right-skewed placement analysis, although a one-sided upper tolerance
limit would better match a direct quality-control decision.

For each of 1000 paired splits, all four methods use the same 100 training
batches and the same 87 held-out batches. TCSP uses the retained count
\(k=98\) from the adaptive scan calibration for \(n=100\), \(c=0.90\), and
tolerance confidence \(0.95\). MTI-ECM uses the frozen selected calibration
policy from the independent simulation validation, with direct
Dirichlet-process content-probability threshold \(0.995\). The fitted content
and tilt are selected by the prespecified MTI-ECM rule within each training
sample; the rule itself is not tuned on these pharmaceutical responses. The
held-out empirical summaries describe stability within the finite historical
data set, with repeated-sampling tolerance assessment left to the simulation
study.

\begin{table}[H]
\centering
\caption{\textbf{Additional weight-RSD application.} Intervals are fit on
100 code-23 batches and evaluated on the 87 held-out batches over 1000 paired
splits. Width and held-out content summaries report medians with empirical
2.5\%--97.5\% ranges across splits.}
\label{rqr:tab:supp-pharma-rsd}
\TableStyle
\begin{tabularx}{\textwidth}{@{}l>{\raggedright\arraybackslash}X>{\raggedright\arraybackslash}X>{\raggedright\arraybackslash}Xr@{}}
\toprule
Method & Median interval & Width median [95\% range] & Held-out content median [95\% range] & Omitted below/above\\
\midrule
TCSP & [0.600, 2.610] & 2.000 [1.527, 2.370] & 96.6 [89.7, 100.0] & 0.0 / 2.3 \\
MTI-ECM & [0.600, 2.610] & 2.000 [1.527, 2.370] & 96.6 [89.7, 100.0] & 0.0 / 2.3 \\
Young--Mathew & [0.700, 2.970] & 2.260 [1.572, 2.589] & 95.4 [87.4, 98.9] & 3.4 / 1.1 \\
Wilks & [0.600, 4.050] & 3.370 [1.818, 3.450] & 98.9 [93.1, 100.0] & 0.0 / 0.0 \\
\bottomrule
\end{tabularx}

\end{table}

Held-out content was also summarized within production month, API batch, and
batch-order quartile whenever a held-out group contained at least three
observations. These checks are descriptive because the splits are drawn from a
finite historical data set rather than from a controlled independent
experiment. Across methods, the median within-split content range is
25.0--33.3 percentage points by API batch, 4.5--10.4 points by batch-order
quartile, and 14.3--33.3 points by production month. The corresponding 97.5\%
summaries reach 100.0, 27.3, and 80.0 points. These large groupwise ranges warn
against interpreting the random-split exercise as evidence under batch or
calendar dependence.

% migration-id: MIG-CH05-05
\section{Computational derivation details}
\subsection{Notation and Targets}
\label{mti:sec:supp-targets}

For an observation \(y_i\) and raw roots \(\eta_{1i}\) and \(\eta_{2i}\),
define
\[
e_i=(y_i-\eta_{1i})(y_i-\eta_{2i}),\qquad
\rho_c(r)=r\{c-\ind(r<0)\}.
\]
The MPI loss is \(\rho_c(e_i)\). With fixed tilt \(\delta_i\), the MTI loss is
\[
\rho_c(e_i)-c\delta_i(\eta_{1i}+\eta_{2i}-2y_i).
\]
The cumulative loss over observed indices \(I_{\mathrm{obs}}\) is denoted by
\(L_{c,\delta}(\theta)\). The fixed-target generalized posterior is
\[
\pi(\theta\mid D)
\propto
\exp\{-\omega_{\mathrm R}L_{c,\delta}(\theta)\}\pi_0(\theta).
\]
The symbols \(c\), \(\delta\), and \(\omega_{\mathrm R}\) have different roles:
content, retained-mean tilt, and learning rate, respectively.

\subsection{Root Labels and Endpoint Summaries}
\label{mti:sec:supp-labels}

The product-residual loss is invariant to exchanging the two raw roots. With
exchangeable root priors, the unordered pair is identified but the labels
``root 1'' and ``root 2'' are not. Endpoint summaries are therefore computed
after sorting the roots at the evaluation points:
\[
L_i=\min(\eta_{1i},\eta_{2i}),\qquad
U_i=\max(\eta_{1i},\eta_{2i}),\qquad
W_i=U_i-L_i.
\]

For static linear roots, define
\[
\beta_m=\frac{\beta_1+\beta_2}{2},\qquad
\beta_s=\frac{\beta_2-\beta_1}{2}.
\]
Then \(M(x)=x^\top\beta_m\), \(S(x)=x^\top\beta_s\), and
\[
L(x)=M(x)-|S(x)|,\qquad
U(x)=M(x)+|S(x)|,\qquad
W(x)=2|S(x)|.
\]
The midpoint coefficient \(\beta_m\) is root-label invariant. The signed gap
coefficient \(\beta_s\) changes sign under a root swap. Lower- and
upper-endpoint coefficient summaries require a declared design region where the
two roots remain ordered.

\subsection{Pseudo-AL Augmentation}
\label{mti:sec:supp-pseudo-al}

Let
\[
\sigma_{\mathrm R}=\omega_{\mathrm R}^{-1},\qquad
\xi_c=\frac{1-2c}{c(1-c)},\qquad
\phi_c=\frac{2}{c(1-c)}.
\]
The single-observation loss factor satisfies
\begin{align}
&\frac{c(1-c)}{\sigma_{\mathrm R}}
\exp\{-\rho_c(e_i)/\sigma_{\mathrm R}\}\notag\\
&\quad=\int_0^\infty
\Normal(e_i\mid \xi_cv_i,\phi_c\sigma_{\mathrm R}v_i)
\Exp(v_i\mid\mathrm{rate}=1/\sigma_{\mathrm R})\,dv_i .
\end{align}
Thus
\[
v_i\mid\cdot\sim
\GIG\left\{
\frac12,\
\frac{1}{2\sigma_{\mathrm R}c(1-c)},\
\frac{c(1-c)e_i^2}{2\sigma_{\mathrm R}}
\right\}.
\]
The latent scales are part of the computational representation of a
loss-based update. They are not latent response variables.

\subsection{Gaussian Root Updates and Fixed Tilt}
\label{mti:sec:supp-static-updates}

Let \(\mat X_{\mathrm{obs}}\) be the observed design matrix and let
\(\mat W_v=\diag\{(\phi_c\sigma_{\mathrm R}v_i)^{-1}:i\in I_{\mathrm{obs}}\}\).
Conditional on the second root, define
\[
\mat A_1=\diag(\vect y_{\mathrm{obs}}-\vect\eta_{2,\mathrm{obs}})
\mat X_{\mathrm{obs}},
\qquad
\vect z_1=\vect y_{\mathrm{obs}}\circ
(\vect y_{\mathrm{obs}}-\vect\eta_{2,\mathrm{obs}})-\xi_c\vect v.
\]
For Gaussian prior canonical parameters \((\mat Q_{01},\vect h_{01})\),
\[
\mat Q_1=\mat Q_{01}+\mat A_1^\top\mat W_v\mat A_1,\qquad
\vect h_1=\vect h_{01}+\mat A_1^\top\mat W_v\vect z_1.
\]
Hence
\[
\beta_1\mid\cdot\sim
\Normal(\mat Q_1^{-1}\vect h_1,\mat Q_1^{-1}).
\]
The second root is updated by exchanging labels.

For fixed nonzero tilt, let \(\vect\delta_{\mathrm{obs}}\) be fixed over the
observed design. The precision matrix is unchanged, and the canonical vector is
shifted to
\[
\vect h_1^{\mathrm{MTI}}
=\vect h_1+
\omega_{\mathrm R}c\mat X_{\mathrm{obs}}^\top\vect\delta_{\mathrm{obs}},
\]
with the same shift for the second root. Proper Gaussian priors give a finite
normalizing constant in finite dimension because the tilt is linear in the
roots.

\subsection{ECM Mode Calculation}
\label{mti:sec:supp-ecm}

The ECM calculation maximizes the fixed generalized-posterior kernel. It uses
the exact inverse-scale moment
\[
\tau_i=E(v_i^{-1}\mid\cdot)=\frac{1}{c(1-c)|e_i|}
\]
away from \(e_i=0\), with a prespecified numerical floor near zero. Holding
\(\beta_2\) fixed, set
\[
\mat B_1=\diag(\vect y_{\mathrm{obs}}-\vect\eta_{2,\mathrm{obs}})
\mat X_{\mathrm{obs}},
\qquad
\vect r_1=\vect y_{\mathrm{obs}}^{\,2}
-\vect y_{\mathrm{obs}}\circ\vect\eta_{2,\mathrm{obs}}.
\]
With ridge prior precision \(\mat P_1\), the conditional maximization solves
\[
\left\{\mat P_1+
\frac{1}{\phi_c\sigma_{\mathrm R}}\mat B_1^\top\diag(\vect\tau)\mat B_1
\right\}\beta_1
=
\frac{1}{\phi_c\sigma_{\mathrm R}}
\mat B_1^\top\{\diag(\vect\tau)\vect r_1-\xi_c\vect 1\}
+\omega_{\mathrm R}c\mat X_{\mathrm{obs}}^\top\vect\delta_{\mathrm{obs}}.
\]
The second root uses the same equation with labels exchanged. A full ECM cycle
updates the latent-scale moments, solves the two conditional Gaussian systems,
and accepts the result only if the exact penalized objective decreases.

\subsection{Conditionally Gaussian Shrinkage}
\label{mti:sec:supp-shrinkage}

The zero-tilt MPI sampler can use a conditionally Gaussian regularized
horseshoe prior through the Nishimura--Suchard representation
\citep{qdesn-CarvalhoPolsonScott2010HS,qdesn-PiironenVehtari2017RHS,exd-nishimura2023shrunken}.
Conditional on the local, global, and slab-scale states, each root coefficient
block has a Gaussian prior precision and therefore fits directly into the
Gaussian update above. This construction is stated for MPI. Extending nonzero
tilt through random shrinkage scales requires its own propriety and diagnostic
checks.

\subsection{Deterministic Basis Expansions}
\label{mti:sec:supp-features}

Let \(\mat X_{\mathrm F}\) denote a deterministic basis matrix obtained before
the endpoint update is fitted. The static formulas apply with
\(\mat X_{\mathrm{obs}}\) replaced by the observed rows of
\(\mat X_{\mathrm F}\). Echo-state and deep echo-state basis functions
\citep{qdesn-Jaeger2001EchoState,qdesn-GallicchioMicheliPedrelli2018DeepESNDesign} are one
source of such a matrix. The recurrent representation is part of the design
construction; posterior uncertainty in the endpoint model is conditional on the
chosen basis unless the basis-generation mechanism is included in the
model.

\subsection{Dynamic Endpoint-State Updates}
\label{mti:sec:supp-dynamic}

A dynamic endpoint model replaces static coefficients by root-specific state
trajectories,
\[
\eta_{jt}=f_t^\top\theta_{jt},\qquad
\theta_{jt}=G_t\theta_{j,t-1}+\varepsilon_{jt},\qquad j=1,2.
\]
Gaussian initial and evolution distributions give a Gaussian prior over each
root path. Because the product residual contains both endpoint paths, the joint
augmented log density is quartic in the two paths. Conditional on one complete
path, however, the other path has a linear Gaussian observation equation.
Root-specific forward-filtering backward-sampling can therefore be alternated
within the specified fixed-target generalized posterior
\citep{env-goncalves}.

Discount factors and component scales are part of the state prior. Fixed values
or precomputed templates define a fixed target. Adaptive choices are best
reported as sensitivity analyses unless they are embedded in a fully specified
probability model.

\subsection{Diagnostics and Validation Limits}
\label{mti:sec:supp-diagnostics}

MCMC summaries should report multiple-chain convergence diagnostics, effective
sample sizes, Monte Carlo error, root-label handling, and endpoint functional
summaries. ECM summaries should report objective descent, iteration counts,
stopping criteria, final roots, fitted content, tilt, and sensitivity to
starting values. Basis-expansion analyses should report how the basis was
chosen. Dynamic analyses should report state-prior choices, smoothing
diagnostics, and whether discount or scale values were fixed in advance.

These diagnostics assess computation for a fixed endpoint target. They do not
establish tolerance validity or response-predictive calibration. Such claims
require a separate validation design aligned with the intended statistical
statement.
